% ===========================================================================
%  Bagian 5.2 -- eblup_sfh
% ===========================================================================
\section{Model Fay--Herriot spasial: \texttt{eblup\_sfh}}\label{sec:sfh}

\subsection{Setting data dan kapan dipakai}\label{sec:sfh-setting}

Bentuk datanya sama dengan FH (satu baris per domain: $y_d$, $\psi_d$, kovariat
$x_d$) tetapi ditambah satu objek wajib: matriks tetangga $W$ berdimensi
$m\times m$ \emph{termasuk} baris/kolom domain tanpa sampel
(\code{R/eblup\_sfh.R:131-141}). Model ini dipakai bila residual direktor
mengandung korelasi spasial, misalnya karena faktor geografis yang tidak
terukur; ekuivalennya adalah model \emph{spatial error}/SAR pada data area-level
\citep{pratesi2008small, rao2015small}. Karena $W$ dari pengguna, fungsi ini
tidak pernah membangun struktur tetangga sendiri; pembuat tetangga rook/queen
tersedia lewat \fct{sim\_spatial\_weights} (bagian~\ref{sec:utilitas}).

\subsection{Persamaan model}\label{sec:sfh-model}

\begin{equation}
  y_d=x_d'\beta+u_d+e_d,\qquad
  u\sim N\bigl(0,\;\sigma^2 A^{-1}\bigr),\qquad
  e_d\sim N(0,\psi_d),
  \label{eq:sfh-model}
\end{equation}
dengan matriks presisi spasial (kode: \code{eblup\_sfh\_core.cpp:311})
\begin{equation}
  A=(\Imat-\rho \Wmat')(\Imat-\rho \Wmat)
   =\Imat-\rho(\Wmat+\Wmat')+\rho^2\Wmat'\Wmat,
  \qquad \abs{\rho}<1.
  \label{eq:sfh-A}
\end{equation}
Karena $A$ adalah presisi, \eqref{eq:sfh-model} ekuivalen dengan bentuk
autoregresi spasial serentak $u=\rho \Wmat u+\varepsilon$, $\varepsilon\sim
N(0,\sigma^2\Imat)$; dalam arti ini efek spasialnya \emph{SAR/err\-or-compo\-nents},
bukan CAR proper (perbandingan kedua notasi ada pada bagian~\ref{sec:car}).
Kovarians total $\Vm=G+\mathrm{diag}(\psi)$ dengan
$G=\sigma^2 A^{-1}$ (baris 318, 401). Turunan yang dibutuhkan skoring
(derifasi lengkap di lampiran~\ref{sec:derivasi}) adalah
\begin{equation}
  \frac{\partial A}{\partial\rho}=-(\Wmat+\Wmat')+2\rho\,\Wmat'\Wmat,
  \qquad
  \frac{\partial \Vm}{\partial\rho}
  =-\sigma^2 A^{-1}\frac{\partial A}{\partial\rho}A^{-1},
  \label{eq:sfh-deriv}
\end{equation}
persis kode \code{eblup\_sfh\_core.cpp:315-316} (matriks $W'W$ dan $W+W'$
di-pra-hitung sekali di luar loop, baris 283-284).

\subsection{Prediktor titik}\label{sec:sfh-prediktor}

Dengan $Q=(\Xmat'\Vm^{-1}\Xmat)^{-1}$ dan
$\betah=(\Xmat'\Vm^{-1}\Xmat)^{-1}\Xmat'\Vm^{-1}y$ (baris 407-414),
prediktor adalah EBLUP matriks persis \eqref{eq:blup}:
\begin{equation}
  \hat\theta=\Xmat\betah+G\,\Vm^{-1}\bigl(y-\Xmat\betah\bigr)
  \quad(\texttt{eblup\_sfh\_core.cpp:462}),
  \qquad
  \hat u=G\Vm^{-1}(y-\Xmat\betah).
  \label{eq:sfh-pred}
\end{equation}
Pada domain tanpa sampel prediktor berubah menjadi \emph{kriging simpel}
(baris 576-591): dengan $G$ dihitung dari $W$ penuh,
\begin{equation}
  \hat u_{ns}=G_{rs}\,\Vm^{-1}\bigl(y_s-\Xmat_s\betah\bigr),
  \qquad
  \hat\theta_{ns}=x_{ns}'\betah+\hat u_{ns}.
  \label{eq:sfh-kriging}
\end{equation}
\implnote{Baris 580-587 memakai sub-blok dari $A_{\text{full}}^{-1}$, sedangkan
pass fit memakai $W_{ss}$; keduanya tidak setara, $[A_{\text{full}}^{-1}]_{ss}
\neq[(\Imat-\rho W_{ss})'(\Imat-\rho W_{ss})]^{-1}$ (SFH-8). Selain itu
dokumentasi menyebut ``analytical MSE approximation'' untuk area tak-sampel,
padahal yang dihitung hanya $g_1+g_2$ tanpa $g_3,g_4$ (SFH-7).}

\subsection{Estimasi komponen varians}\label{sec:sfh-varians}

Parameter yang diiterasikan dua: $\theta=(\sigma_u^2,\rho)$, nilai awal
$\sigma_u^{2(0)}=\mathrm{median}(\psi)$ dan $\rho^{(0)}=0.5$
(baris 288-289), toleransi relatif $10^{-4}$, maksimum 100 iterasi
(baris 307). Per iterasi (baris 308-383):
\begin{enumerate}[label=\arabic*.]
  \item bangun $A$, $A^{-1}$ (\code{inv\_sympd}, fallback \code{pinv}),
        $\Vm$, $\Vm^{-1}$, $Q$, $P=\Vm^{-1}-\Vm^{-1}\Xmat Q\Xmat'\Vm^{-1}$;
  \item skor dan informasi --- untuk REML (baris 349-357)
        \begin{equation}
          s_a=-\tfrac12\tr(P D_a)+\tfrac12 y'P D_a P y,
          \qquad
          \mathcal I_{ab}=\tfrac12\tr(P D_a P D_b),
          \label{eq:sfh-score-reml}
        \end{equation}
        dengan $D_1=A^{-1}=\partial\Vm/\partial\sigma^2$ (kode: \code{derSigma}) dan $D_2=\partial\Vm/\partial\rho$; untuk ML
        (baris 338-348) baris pertama skor memakai $\Vm^{-1}$ tetapi
        kuadratiknya tetap $y'PDPy$;
  \item langkah Newton--Fisher $\Delta=\mathcal I^{-1}s$ lewat \code{solve};
        kegagalan memutus loop dengan \code{convergence = FALSE};
  \item batas: $|\rho|<0.999$ (baris 371-374); $\sigma_u^2$ \emph{tidak}
        dibatasi selama iterasi, hanya $\sigma^2=\max(0,\cdot)$ pada hasil
        akhir (baris 394);
  \item konvergensi $\texttt{diff}=\max_a\abs{\Delta\theta_a/(\theta_a+10^{-12})}$
        (baris 376-383).
\end{enumerate}
Galat baku kedua komponen $\sqrt{\mathrm{diag}(\mathcal I^{-1})}$ dihitung
dari $\mathcal I$ berbentuk REML tanpa memeriksa \code{method} (baris 502-510,
647-653), dan \code{loglikelihood/AIC/BIC} selalu memakai bentuk ML
(baris 464-476).

\implnote{Tiga hal pada penutup paragraf di atas terdaftar sebagai SFH-1
(skor ML setengah REML), SFH-2 (SE dari informasi REML pada fit ML), dan
SFH-3 (loglik ML walau REML) di \code{DISCREPANCIES.md}. Batas akhir
$\rho=\pm0.999$ dipetakan ke $\pm1.0$ persis tempat $A$ singular lalu jatuh ke
\code{pinv} (baris 389-399), sedangkan jalur bootstrap mempertahankan
$\pm0.999$ --- dua perilaku berbeda untuk kasus yang sama (SFH-4).}

\subsection{Estimasi MSE}\label{sec:sfh-mse}

Empat suku dihitung setelah loop konvergen (baris 481-543). Dengan
$R=\Xmat-G\Vm^{-1}\Xmat$ (kode: \code{mat R = X - Gb}):
\begin{align}
  g_1 &= \diag\bigl(G-G\Vm^{-1}G\bigr),
    \label{eq:sfh-g1}\\
  g_2 &= \diag\bigl(R\,Q\,R'\bigr),
    \label{eq:sfh-g2}\\
  g_3(i) &= \tr\bigl(L_i\,\Vm L_i'\,\mathcal I^{-1}\bigr),
    \label{eq:sfh-g3}\\
  g_4(i) &= \tfrac12\,D_{ii},
    \label{eq:sfh-g4}
\end{align}
dengan $L_i=(l_{1i}',\,l_{2i}')'$,
$l_1=\Vm^{-1}D_1-\sigma^2\Vm^{-1}D_1\Vm^{-1}D_1$,
$l_2=\Vm^{-1}D_2-\sigma^2\Vm^{-1}D_2\Vm^{-1}D_1$,
$D=\bigl(\Psi V^{-1}D_{12}V^{-1}\Psi\bigr)(\mathcal I^{-1}_{01}+\mathcal
I^{-1}_{10})+\bigl(\Psi V^{-1}D_{22}V^{-1}\Psi\bigr)\mathcal I^{-1}_{11}$,
$\Psi=\mathrm{diag}(\psi)$ (baris 512-541). Lalu
\begin{equation}
  \widehat{\mathrm{MSE}}(\hat\theta_d)=g_1+g_2+2g_3-g_4
  \quad(\texttt{baris 543}),
  \qquad
  \text{dan bila \code{method = "ML":}\ }
  \widehat{\mathrm{MSE}}\mathrel{-}=\mathbf{b}_{\text{ML}}'
  \nabla g_1 .
  \label{eq:sfh-mse}
\end{equation}
Koreksi bias ML (baris 546-571) memakai
$b_a=-\tr(Q\Xmat'\Vm^{-1}D_a\Vm^{-1}\Xmat)$ dan $b_{\text{ML}}
=\tfrac12\mathcal I^{-1}b$, tetapi hanya menurunkan $g_1$ --- suku lain
tidak dikoreksi. Pada area tak-tersampel berlaku
$\mathrm{MSE}=g_1+g_2$ saja (baris 596-599).

\paragraph{Varian bootstrap.} Dua pilihan tambahan lewat
\code{mse\_method}:
\begin{itemize}
  \item \code{"pbmse"}: bootstrap parametrik
        \eqref{eq:pbmse} --- replikat $y^\ast=X\hat\beta+v^\ast+e^\ast$,
        $v^\ast=(\Imat-\rho W)^{-1}u^\ast$, $u^\ast\sim N(0,\hat\sigma^2)$;
        refit Fisher scoring penuh per replikat dalam region OpenMP
        (algoritme~\ref{alg:bootstrap}); skema batch menjamin tepat $B$
        replikat valid (\code{eblup\_sfh\_pbmse.cpp:77-88}).
        Versi bias-koreksi \eqref{eq:pbbc} dihitung bersamaan.
        Kerangka bootstrap untuk EBLUP spasial dirujuk ke
        \citet{molina2009bootstrap}.
  \item \code{"npbmse"}: bootstrap nonparametrik atas residual
        terstandarisasi; error sampling dan efek acak dipisahkan lewat
        eigen-dekomposisi $\Psi_e$ dan $\Psi_u$ (membuang $p$ eigen
        terkecil, sesuai proyeksi desain), lalu indeks residual
        diresample; hanya untuk \code{method = "REML"} --- pemeriksaannya
        berada di C++, bukan di R (\code{eblup\_sfh\_npbmse.cpp:32-34};
        SFH-9).
\end{itemize}
Kedua bootstrap berjalan pada subset area tersampel, lalu hasil fit ulang
penuh ditimpa (baris tersampel) oleh hasil bootstrap
(\code{R/eblup\_sfh.R:155-200}).

\subsection{Pseudo-code}\label{sec:sfh-algo}

\begin{algorithm}[htbp]
\caption{\texttt{eblup\_sfh} --- alur sebenarnya (\texttt{R/eblup\_sfh.R:85-239}, \texttt{src/eblup\_sfh\_core.cpp:210-670})}
\label{alg:sfh}
\KwIn{formula, \code{vardir}, \code{data}, $W$ (wajib), \code{method},
  \code{mse\_method}, \code{B}, \code{n\_threads}, \code{seed}}
\KwOut{\code{df\_eblup} (+ kolom bootstrap bila ada), \code{rho},
  \code{estvarcomp}, \code{goodness}}
\tcp*{\emph{Lapis R}}
\code{method} $\leftarrow$ \texttt{match.arg(toupper(...))} (100-101)\;
\textbf{gagal} bila \code{W} kosong atau bukan $n\times n$ (131-141)\;
\code{y} $\leftarrow$ respons (\code{na.pass}, 109); $I_s\leftarrow$ \texttt{!is.na(y)} (144)\;
\uIf{\code{mse\_method} $\in$ \{pbmse, npbmse\}}{
  subset $X_s,y_s,\psi_s,W_{ss}$ (155-158)\;
  $r_b\leftarrow$ \texttt{.seblup\_pbmse}/\texttt{.seblup\_npbmse} (161, 167)\;
  bila ada area tak-sampel: \texttt{.seblup\_core} penuh, timpa baris $I_s$
  dengan hasil bootstrap, kolom bootstrap area tak-sampel $\leftarrow$ \texttt{NA} (179-200)
}
\uElse{\texttt{.seblup\_core($\Xmat,y,\psi,W$, method, maxiter, precision)} penuh (205-213)}\;
\caption*{}
\tcp*{\emph{Lapis C++}}
\textbf{gagal} bila $W$ bukan bujur sangkar vs \code{nrow(X)} (223-225)\;
$I_{ns}\leftarrow$ \texttt{find\_nonfinite}(y); $W\leftarrow W[I_s,I_s]$ (240-257)\;
pra-hitung $W',\Imat,W'W,W+W'$ (279-284)\;
$\sigma_u^2\leftarrow\mathrm{median}(\psi)$;\quad $\rho\leftarrow0.5$ (286-289)\;
\While{\texttt{diff} $>$ \code{precision} \textbf{dan} $k<$ \code{maxiter}}{
  $A\leftarrow(\Imat-\rho W')'(\Imat-\rho W)$;\ $A^{-1}\leftarrow$ \texttt{inv\_sympd}, fallback \texttt{pinv} (308-313)\;
  $\partial A/\partial\rho$, $\partial V/\partial\rho\leftarrow$ \eqref{eq:sfh-deriv} (315-316)\;
  $V\leftarrow\sigma^2A^{-1}+\diag(\psi)$;\ $V^{-1}\leftarrow$ \texttt{inv\_sympd}, fallback \texttt{pinv} (318-320)\;
  $Q\leftarrow(\Xmat'V^{-1}\Xmat)^{-1}$;\ $P\leftarrow V^{-1}-V^{-1}\Xmat Q\Xmat'V^{-1}$ (323-331)\;
  \uIf{ML}{$s_a\leftarrow-\tfrac12\tr(V^{-1}D_a)+\tfrac12 y'PDPy$;
    $\mathcal I_{ab}\leftarrow\tfrac12\tr(V^{-1}D_aV^{-1}D_b)$ (338-348)}
  \uElse{$s_a,\mathcal I_{ab}\leftarrow$ \eqref{eq:sfh-score-reml} (349-357)}\;
  $\Delta\leftarrow\mathcal I^{-1}s$; bila gagal: \texttt{convergence = FALSE}, hentikan (365-369)\;
  $\rho\leftarrow\mathrm{clamp}(\rho+\Delta_\rho,\,-0.999,\,0.999)$;\quad
  $\sigma_u^2\leftarrow\sigma_u^2+\Delta_{\sigma}$ (371-374)\;
  $\texttt{diff}\leftarrow\max_a\abs{\Delta\theta_a/(\theta_a+10^{-12})}$; non-finite $\to$ gagal (376-383)\;
}
$\rho_{\text{fix}}\leftarrow\rho$; bila $=\pm0.999\to\pm1.0$;\quad $\sigma^2\leftarrow\max(0,\sigma_u^2)$ (389-394)\;
bangun ulang $A^{-1},G,V,V^{-1},Q,\betah$ (396-416)\;
galat baku, $z$, $p$ (normal), $\text{resid}\leftarrow y-\Xmat\betah$ (451-460),
$\hat\theta\leftarrow$ \eqref{eq:sfh-pred} (462)\;
loglik ML, AIC, BIC (464-476)\;
$g_1,g_2,g_3,g_4\leftarrow$ \eqref{eq:sfh-g1}--\eqref{eq:sfh-g4};
  $\mathrm{MSE}\leftarrow g_1+g_2+2g_3-g_4$ (481-543)\;
\If{\code{method} $=$ ML}{koreksi $\mathrm{MSE}\mathrel{-}=\mathbf b_{\text{ML}}'\nabla g_1$ (546-571)}\;
\For{$d\in I_{ns}$}{kriging \eqref{eq:sfh-kriging};\ $\mathrm{MSE}\leftarrow g_1+g_2$ (576-604)}\;
$\mathrm{RSE}\leftarrow100\sqrt{\mathrm{MSE}}/\abs{\hat\theta}$; susun keluaran (621-669)\;
\end{algorithm}

\subsection{Signature, argumen, objek keluaran, dan contoh}\label{sec:sfh-api}

\begin{lstlisting}[language=R,caption={Signature \texttt{eblup\_sfh} (\texttt{R/eblup\_sfh.R:85-99}).}]
eblup_sfh(formula, vardir, domain = NULL, data,
          method = c("REML", "ML"),
          mse_method = c("analytical", "pbmse", "npbmse"),
          W = NULL, B = 100, n_threads = 1, seed = -1,
          maxiter = 100, precision = 1e-4, print_result = TRUE)
\end{lstlisting}

\begin{itemize}
  \item \code{W} --- matriks tetangga $n_{\text{total}}\times n_{\text{total}}$
        (row-standar/sesuai skala yang diinginkan pengguna); \emph{wajib},
        tidak ada default.
  \item \code{mse\_method} --- \code{"analytical"} (default), \code{"pbmse"},
        \code{"npbmse"}.
  \item \code{B} --- jumlah replikat bootstrap (default 100);
        \code{n\_threads} --- $>0$ memakai \texttt{omp\_set\_num\_threads},
        $\le0$ memakai semua core; \code{seed = -1} hanya melewatkan
        penyetelan ulang benih (\code{if (seed >= 0) set\_r\_seed(seed)},
        \code{eblup\_sfh\_pbmse.cpp:29}), sehingga draw bootstrap tetap
        memakai RNG R yang sedang berjalan walau dokumentasi menyebutnya
        ``leaves the current R RNG state unchanged''.
  \item \code{method} --- REML (default) atau ML.
\end{itemize}

Objek keluaran: \code{"fastsae"} dengan \code{estcoef}
(\code{beta, std.error, stderr\_beta, zvalue, pvalue}), \code{random\_effect\_var},
\code{rho}, \code{estvarcomp} (2 baris: \code{sigma2\_u}, \code{rho} ---
keduanya dari informasi REML), \code{goodness}
(\code{loglikelihood, AIC, BIC}), \code{df\_eblup} (\code{y, vardir, eblup,
random\_effect, mse, rse} + kolom bootstrap), \code{W}, \code{n\_iter},
\code{convergence}, \code{model = "SFH"}, \code{level = "area"}.
Kolom galat baku namanya \code{zvalue}, bukan \code{tvalue} seperti tertulis
di dokumentasi (SFH-6).

\begin{lstlisting}[language=R,caption={Contoh yang dapat dijalankan (dari \texttt{man/eblup\_sfh.Rd}).}]
library(fastsae)
m1 <- eblup_sfh(y ~ x1 + x2 + x3, data = mys, vardir = ~vardir,
                W = mys_proxmat)
m1$rho                       # autoregressive spatial parameter
m1$estvarcomp                # sigma2_u and rho with std.error

m2 <- eblup_sfh(y ~ x1 + x2 + x3, data = mys, vardir = ~vardir,
                W = mys_proxmat, mse_method = "pbmse", B = 50)
\end{lstlisting}

\subsection{Referensi kunci}\label{sec:sfh-ref}

\citet{pratesi2008small} mengembangkan FH spasial; \citet{rao2015small}
menyajikan kerangka EBLUP yang dipakai untuk semua suku $g$;
blok \code{@references} fungsi ini hanya memuat \citet{rao2015small}
(\code{R/eblup\_sfh.R:5-8}). Referensi pembanding implementasi:
\pkg{sae} \citep{molina2015sae}.
